\section{A Pre-Registered Program}
\label{sec:program}

None of the following has run. Each experiment states its success criterion and
falsifier before any result, as every measurement in this work has. They are ordered
so that each one measures the layer it targets, not a gap an earlier experiment could
have closed.

\paragraph{E1: agent-written adherence proofs, across languages.}
This experiment asks whether an agent can produce a certificate (\S\ref{sec:formal})
that a kernel adheres to its spec. It also asks in which formal languages, for which
kernels, at what cost, and how often an accepted proof certifies the wrong thing.
\begin{itemize}
  \item \emph{Targets.} There are \pilotTargets{} of them, and the agent is not told which are wrong.
    Three are correct: the Triton softmax and layernorm, and an \texttt{exp2} softmax that is
    correct only because $2^{x\log_2 e} = e^x$. Four are wrong: $\softmax(2x)$, the
    constant fill, rolled rows, and a kernel that drops each row's tail. For each, the
    agent proves adherence or refutes it.
  \item \emph{Languages.} \pilotLanguages{} in the pilot, spanning three verifier families:
    Lean~4 with Mathlib \citep{demoura2021lean4,mathlib2020}, both with and without VeriTile
    \citep{veritile2026}, plus Rocq and Isabelle/HOL; the SMT-driven Dafny, Verus, F* and
    Why3; and floating-point routes through Rocq with Flocq and Gappa and through Why3's
    SMT floating-point theories. Other languages join through adapters. A language that
    cannot state the spec at all is recorded as a result.
  \item \emph{Protocol.} Each (language, kernel, replicate) attempt is a fresh Claude
    sub-agent with no budget cap: \pilotSessions{} sessions at \pilotReplicates{}
    replicates. Cost is recorded for every session. The agent sees the six rules and none
    of the data, the held-out practice of \citet{ye2025verina} and \citet{thakur2025clever}.
\end{itemize}
\emph{Measured:} the certified rate on correct kernels (pass@$k$); \emph{false
certificates} on wrong kernels; the \emph{gaming profile}, meaning attempts that C1 accepts
but C2--C6 reject, broken down by check; spec completeness on the hidden set against the
human-written bundle's; refutations; and cost. \emph{Success} requires zero false
certificates. Any single one is a hole in the harness and stops the study until it is
closed. \emph{Predicted:} Dafny reaches the highest C1 rate and Lean the lowest, as
\citet{bursuc2025vericoding} found for general programs. We also predict that in at least
one language the agent writes a spec made of laws, which C3 rejects.

\paragraph{E2: complete the bundle.}
Proposition~\ref{prop:softmax}(a) and (b) are mechanized; mechanize (c). Then add the log-ratio and affine-tie laws to
the gate with tolerances, and cover the inputs the real call site sends, including
$-\infty$ masked scores. \emph{Success:} every family the proof says the old bundle admits
is admitted by the old gate and rejected by the new one. The correct kernel and the
positive control stay as they are, and no new false positive appears on the
\loopCandidates{} loop candidates or the \fpGroups{} device groups. \emph{Falsified} by any
new false positive. Whether fast-$\exp$ is admitted is decided by a stated accuracy
budget.

\paragraph{E3: red-team the gate.}
Run a loop whose objective is kernel runtime, with no list of forbidden moves. Agents
see only the gate's verdicts, and a hidden checker (the reference arm plus masked
rows) re-scores every admitted kernel. Two conditions differ only in the bundle, current
or completed. \emph{Measured:} the rate of wrong kernels admitted, time to the first
hack, the best correct speedup, and gate latency. \emph{Predicted:} the current bundle is
hacked, with the constant fill as the obvious target, and the completed bundle is not. The
best correct speedups match.

\paragraph{E4: formal feedback for self-improvement.}
Repeat E3 with the completed bundle in both conditions, with and without formal results
returned to the agent: failed checks, unsolved goals and refuting inputs. This tests
whether proof feedback improves attempts beyond what a better gate alone gives.
Verifier feedback is known to help in verified code generation
\citep{ye2025verina,aggarwal2025alphaverus}, but not yet for numerical kernels.

\paragraph{Cost, throughout.}
We log the time and tokens each formalization takes. If formalizing a task costs more
than writing a trusted reference, the declarative arm's ``no reference needed''
argument weakens, and that is also a result.
